\begin{lemma}[Negative scaling of infinity-harmonic functions]
Let $E$ be a finite-dimensional real inner-product space, let $\Omega\subseteq E$,
let $u:E\to\R$ be infinity-harmonic on $\Omega$ in the sense of
\noderef{IsInfHarmonicOn}, and let $t<0$. Then the function $x\mapsto tu(x)$ is
again infinity-harmonic on $\Omega$ in the sense of \noderef{IsInfHarmonicOn}.
\end{lemma}
\begin{proof}
By \noderef{IsInfHarmonicOn} the hypothesis says that $u$ is both an infinity
subsolution on $\Omega$ in the sense of \noderef{IsInfSubsolutionOn} and an
infinity supersolution on $\Omega$ in the sense of
\noderef{IsInfSupersolutionOn}, and the conclusion to be proved is the
conjunction of the same two properties for $tu$. Unwinding
\noderef{IsInfSubsolutionOn} and \noderef{IsInfSupersolutionOn}, each of those
two properties consists of three clauses: that $\Omega$ be open, that the
function be continuous on $\Omega$, and a one-sided inequality at every $C^2$
test. The first clause is part of the hypothesis and is unchanged. For the
second, $u$ is continuous on $\Omega$ by the hypothesis, hence so is the
multiple $tu$ of it by a constant. There remain the two test inequalities.

\emph{The scaling identity.} Let $s\in\R$, let $z\in E$ and let $\varphi$ be a
$C^2$ function near $z$. Then
\begin{equation}\label{eq:ihns-scale}
  \Dinf(s\varphi)(z)=s^3\,\Dinf\varphi(z).
\end{equation}
Indeed $\nabla(s\varphi)(z)=s\nabla\varphi(z)$ and
$D^2(s\varphi)(z)=sD^2\varphi(z)$, these being the first and second derivatives
of a constant multiple; so, by the definition of $\Dinf$ in \noderef{InfLap} and
the bilinearity of the second derivative,
\[
  \Dinf(s\varphi)(z)=D^2(s\varphi)(z)\bigl[\nabla(s\varphi)(z)\bigr]
   =s\,D^2\varphi(z)\bigl(s\nabla\varphi(z),s\nabla\varphi(z)\bigr)
   =s^3D^2\varphi(z)\bigl[\nabla\varphi(z)\bigr]
   =s^3\Dinf\varphi(z).
\]

\emph{$tu$ is a subsolution.} Let $z\in\Omega$ and let $\varphi$ be $C^2$ at $z$
with $tu-\varphi$ having a local maximum on $\Omega$ at $z$, i.e.
$tu(x)-\varphi(x)\le tu(z)-\varphi(z)$ for all $x$ in $\Omega$ near $z$.
Multiplying this inequality by $t^{-1}<0$ reverses it:
\[
  u(z)-t^{-1}\varphi(z)\le u(x)-t^{-1}\varphi(x)
\]
for all such $x$, since $t^{-1}(tu(w)-\varphi(w))=u(w)-t^{-1}\varphi(w)$ for
every $w$. Thus $t^{-1}\varphi$, which is again $C^2$ at $z$, is a lower $C^2$
test for $u$ at $z$: $u-t^{-1}\varphi$ has a local minimum on $\Omega$ at $z$.
As $u$ is an infinity supersolution on $\Omega$,
\noderef{IsInfSupersolutionOn} gives $\Dinf(t^{-1}\varphi)(z)\le0$. By
\eqref{eq:ihns-scale} with $s=t^{-1}$ this reads $t^{-3}\Dinf\varphi(z)\le0$.
Since $t<0$ we have $t^{-1}<0$ and hence $t^{-3}=t^{-1}(t^{-1})^2<0$; dividing
by $t^{-3}$ reverses the inequality and yields $\Dinf\varphi(z)\ge0$. As $z$ and
$\varphi$ were arbitrary, $tu$ is an infinity subsolution on $\Omega$ in the
sense of \noderef{IsInfSubsolutionOn}.

\emph{$tu$ is a supersolution.} Symmetrically, let $z\in\Omega$ and let
$\varphi$ be $C^2$ at $z$ with $tu-\varphi$ having a local minimum on $\Omega$
at $z$. Multiplying $tu(z)-\varphi(z)\le tu(x)-\varphi(x)$ by $t^{-1}<0$ gives
$u(x)-t^{-1}\varphi(x)\le u(z)-t^{-1}\varphi(z)$ for $x$ in $\Omega$ near $z$,
so $t^{-1}\varphi$ is an upper $C^2$ test for $u$ at $z$. As $u$ is an infinity
subsolution on $\Omega$, \noderef{IsInfSubsolutionOn} gives
$0\le\Dinf(t^{-1}\varphi)(z)=t^{-3}\Dinf\varphi(z)$ by
\eqref{eq:ihns-scale}, and dividing by $t^{-3}<0$ yields
$\Dinf\varphi(z)\le0$. Hence $tu$ is an infinity supersolution on $\Omega$ in
the sense of \noderef{IsInfSupersolutionOn}.

The two displayed properties together are exactly infinity-harmonicity of $tu$
on $\Omega$ in the sense of \noderef{IsInfHarmonicOn}.
\end{proof}
